\section{Which open questions should be pursued next?}

\subsection{First priority: the six-term singular-modulus problem}

The prior classification \eqref{eq:prior-classification} gives a concrete next question:
\[
 \text{If }h(\Delta)\ge6,\quad\text{must }s(\Delta)\ge7?
\]
The two exact powered families prove partial cases of the corresponding support-seven phenomenon. They do not settle the original six-term question.

An effective next stage would start with the remaining low-degree orders and certify all exponent residue patterns using the stronger blockwise test. The required output is finite and precise: exact class-polynomial coefficients, scalar congruences modulo $p^{s+1}$, and a covering argument showing that every candidate six-exponent pattern is rejected by at least one rank test. A successful global statement would additionally need a large-discriminant reduction, as the five-term report supplied. The finite and infinite parts should be proved separately and then joined.

This target has three advantages. It is already connected to a published question about rational sparse equations for singular moduli;\cite{blp} it has exact arithmetic certificates; and the new local theorem changes which small-prime tests are available. The generic fewnomial literature provides context, while the CM conjugates and the finite residue covering supply the special arithmetic content.\cite{fewnomials}

\subsection{Second priority: an arithmetic family of constrained twists}

For the Selmer route, the next useful deliverable is not another isolated torsion point. It is a family that simultaneously has:

\begin{enumerate}
\item an explicit ray-group presentation and required-triviality subgroup $D$;
\item a proved nonvanishing or bounded-valuation witness on the compatible character subgroup;
\item controlled local orders so the first required powers remain ramified;
\item a quantitative support-prime construction meeting the actual Frobenius conditions; and
\item a bound for the associated conductor and coefficient-field ramification.
\end{enumerate}

Theorems~\ref{thm:extension}--\ref{thm:bounded-conductor} solve parts of this interface and show exactly why the parts do not collapse into one assertion. The $151$ example is a complete elementary benchmark. A harder benchmark would replace its single class relation by the tied cocycle relations of one explicitly chosen Selmer datum.

The auxiliary primes in the source are selected in prescribed Frobenius neighborhoods. One must identify the governing extension before deciding which zero theorem or effective Chebotarev bound is available. Rational congruences only control abelian information. This is where a purported conductor bound based only on Dirichlet GRH can fail.

\subsection{Nonabelian GRH: aim at one arithmetic inequality}

Earlier reports already reduce several nonabelian zeta questions to explicit one-sided prime-sum estimates. The cubic field $\Q(\alpha)$ with $\alpha^3-\alpha-1=0$, and its $S_3$ splitting field, are natural small benchmarks; a class-sensitive Hecke character over $\Q(\sqrt{-23})$ occurs there. The new character calculation concerns auxiliary ray characters of that quadratic field. It does not prove a zero theorem for its class-character $L$-function.

Further progress should be measured by a new arithmetic input: an integral lattice with the right Frobenius-sensitive divisibility, an independent family, a cheap completion, or a determinant lower bound with a controlled error. Proposition~\ref{prop:rank} allows these inputs to be budgeted before an ambitious construction is attempted.

The repository endpoint cannot simply be applied to the nonabelian factor. Nor does finding anticyclotomic characters imply that the theta lifts or Selmer hypotheses needed for its determinant comparison have been constructed.

\subsection{\texorpdfstring{$K$}{K}-groups, Vandiver, and BSD}

For the $K$-theory route, the most valuable next result is the comparison between the source's custom odd-cocycle module and the conventional cohomological object, with all local conditions and exceptional primes recorded. After that, a finite Hecke determinant certificate may become a proof of a genuine arithmetic statement. Before it, the computation is a candidate test.

For Vandiver-type questions, parity, field variation, and the distinction between a normalized length and actual vanishing remain material. Passing to a more ramified coefficient field can make normalized lengths small without killing a module. The global-order increase of Section~\ref{sec:characters} should therefore be recorded explicitly, not hidden in a change of coefficient ring.

For BSD, the general twisting theorem is not a theorem about arbitrary elliptic-curve Selmer groups. The scalar-inertia and conjugate-dual conditions must be established for the actual representation. Anticyclotomic twisting can preserve a conjugate-dual obstruction rather than remove it. A direct substitution remains unsupported.

\subsection{Research verdict}

There is real transferable mathematics here. The completed CM classification and the new powered families show a tractable direction in which arithmetic divisibility and determinants give exact results beyond familiar RH consequences. The character calculation gives another concrete output and corrects a tempting but false identification between local torsion data and global twists of the same order.

The path to a new broad zero theorem remains open at its arithmetic core. The practical way forward is to prove a missing construction or estimate for one chosen target, with a numerical error budget such as \eqref{eq:rank-cost}, and to stop counting equivalent analytic criteria as resolutions.
